\documentclass[12pt]{report}

\usepackage[a4paper,margin=1in]{geometry}
\usepackage[T1]{fontenc}
\usepackage[utf8]{inputenc}
\usepackage{lmodern}
\usepackage{graphicx}
\usepackage{booktabs}
\usepackage{tabularx}
\usepackage{array}
\usepackage{ragged2e}
\usepackage{enumitem}
\usepackage{listings}
\usepackage{xcolor}
\usepackage{float}
\usepackage[hidelinks]{hyperref}
\usepackage{caption}

\captionsetup{font=small}
\setlist[itemize]{leftmargin=*,itemsep=0.35em}
\setlist[enumerate]{leftmargin=*,itemsep=0.35em}
\setlength{\parskip}{0.6em}
\setlength{\parindent}{0pt}
\renewcommand{\arraystretch}{1.25}

\definecolor{codebg}{RGB}{248,248,248}
\definecolor{placeholdergray}{RGB}{110,110,110}

\lstdefinestyle{reportcode}{
  basicstyle=\ttfamily\small,
  backgroundcolor=\color{codebg},
  frame=single,
  breaklines=true,
  columns=fullflexible,
  keepspaces=true
}

\newcolumntype{Y}{>{\RaggedRight\arraybackslash}X}
\newcommand{\placeholder}[1]{\textcolor{placeholdergray}{\textit{#1}}}
\newcommand{\evidencebox}[1]{
  \begin{figure}[H]
    \centering
    \fbox{\rule{0pt}{2.1in}\rule{0.88\textwidth}{0pt}}
    \caption{#1}
  \end{figure}
}

\newcommand{\lessonchapter}[3]{
  \chapter{Lesson #1: #2}

  \section{Part 1: Goal and Vulnerability Summary}
  \placeholder{State the vulnerability type, main affected component, security impact, and the expected safe behavior for #2.}

  \section{Part 2: Why This Works / Root Cause}
  \placeholder{Explain the security mistake that makes this #3 issue possible. Keep this conceptual and focused on root cause rather than symptoms.}

  \section{Part 3: Environment and Setup}
  \begin{itemize}
    \item DVSA Website URL / API endpoint: \placeholder{Fill in the exact endpoint or page used.}
    \item AWS services involved: \placeholder{For example API Gateway, Lambda, Cognito, S3, DynamoDB, SQS, IAM.}
    \item Relevant function, role, or log group: \placeholder{Name the exact resources used in this lesson.}
    \item User or test context: \placeholder{Attacker, victim, admin, or anonymous context as applicable.}
    \item Tools used: \placeholder{Browser DevTools, curl, jq, AWS Console, AWS CLI, Policy Simulator, CloudWatch, etc.}
  \end{itemize}

  \section{Part 4: Reproduction Steps}
  \begin{enumerate}
    \item \placeholder{Describe the baseline or normal behavior first.}
    \item \placeholder{List the exact steps that trigger the vulnerability.}
    \item \placeholder{Include the request, timing, role, or payload details needed for reproducibility.}
    \item \placeholder{Record the observable result that proves the issue.}
  \end{enumerate}

  \section{Part 5: Evidence and Proof}
  \begin{itemize}
    \item \placeholder{Attach the strongest pre-fix evidence for this lesson.}
    \item \placeholder{Use only evidence that directly proves the claim.}
    \item \placeholder{Redact real tokens, secrets, and temporary credentials before submission.}
  \end{itemize}
  \evidencebox{Replace with pre-fix evidence for Lesson #1 (#2).}

  \section{Part 6: Fix Strategy / Probable Mitigation}
  \begin{itemize}
    \item Fix location: \placeholder{Identify the backend, IAM, cloud configuration, or workflow layer that must change.}
    \item Security property restored: \placeholder{For example authentication integrity, authorization, least privilege, or safe input handling.}
    \item Why this addresses root cause: \placeholder{Explain why the fix is real and not cosmetic.}
  \end{itemize}

  \section{Part 7: Code / Config Changes}
  \placeholder{Paste the real code diff, policy change, CloudFormation change, or configuration update below.}

  \begin{figure}[H]
    \centering
    \fbox{
      \parbox{0.88\textwidth}{
        \vspace{0.4em}
        \placeholder{Replace this box with the actual diff, code snippet, policy JSON, or configuration change.}
        \vspace{2.0in}
      }
    }
    \caption{Replace with the concrete code or configuration change for Lesson #1.}
  \end{figure}

  \section{Part 8: Verification After Fix}
  \begin{enumerate}
    \item \placeholder{Repeat the same test that proved the vulnerability.}
    \item \placeholder{Show that the exploit no longer succeeds.}
    \item \placeholder{Also show that normal intended behavior still works.}
  \end{enumerate}
  \evidencebox{Replace with post-fix verification evidence for Lesson #1 (#2).}

  \section{Part 9: Structured Operation and Security Analysis}
  \subsection*{Table A}
  \begin{table}[H]
    \small
    \centering
    \begin{tabularx}{\textwidth}{|Y|Y|Y|Y|}
      \hline
      \textbf{Intended rule(s)} & \textbf{Artifacts used to infer the rule} & \textbf{Normal evidence} & \textbf{Exploit evidence} \\
      \hline
      \placeholder{State the intended security or workflow rule.} &
      \placeholder{Identify the source used to infer the rule: UI flow, code path, IAM policy, logs, helper guide, etc.} &
      \placeholder{Show normal behavior when the rule is respected.} &
      \placeholder{Show behavior proving the rule was violated.} \\
      \hline
    \end{tabularx}
  \end{table}

  \subsection*{Table B}
  \begin{table}[H]
    \footnotesize
    \centering
    \begin{tabularx}{\textwidth}{|Y|Y|Y|Y|Y|}
      \hline
      \textbf{Why the exploit is a deviation} & \textbf{Deviation class} & \textbf{Fix applied} & \textbf{Post-fix verification} & \textbf{Optional latency / timing note} \\
      \hline
      \placeholder{Explain why the observed result violates the intended rule.} &
      \placeholder{Examples: authentication failure, access control failure, unsafe input handling, logic flaw, DoS, excessive privilege.} &
      \placeholder{Summarize the real corrective action.} &
      \placeholder{Describe the evidence that the deviation is closed.} &
      \placeholder{Use this only when timing or latency matters.} \\
      \hline
    \end{tabularx}
  \end{table}

  \section{Part 10: Takeaway / Lessons Learned}
  \placeholder{End with a short security lesson that ties the #3 weakness back to serverless design, AWS service interaction, and secure engineering practice.}

  \clearpage
}

% Fill in these fields before compiling the final report.
\newcommand{\CourseCode}{ICS-344}
\newcommand{\CourseTitle}{Information Security}
\newcommand{\CourseSection}{Section 1 \& 2}
\newcommand{\CourseTerm}{252}
\newcommand{\ProjectTitle}{DVSA Vulnerability Discovery and Remediation}
\newcommand{\StudentNames}{Husain Al Muallim, Ali Al Qatari}
\newcommand{\StudentIDs}{202163730, 202279780}
\newcommand{\DVSAWebsiteURL}{\placeholder{https://your-dvsa-site.example}}
\newcommand{\AWSRegion}{us-east-1}
\newcommand{\SubmissionDate}{\placeholder{Submission date}}

\begin{document}

\hypersetup{pageanchor=false}
\begin{titlepage}
  \centering
  {\Large King Fahd University of Petroleum and Minerals\par}
  \vspace{1.2cm}
  {\Large \CourseCode\ - \CourseTitle\par}
  {\large \CourseSection\par}
  \vspace{1.4cm}
  {\huge \bfseries \ProjectTitle\par}
  \vspace{1.4cm}

  \begin{tabular}{rl}
    Term: & \CourseTerm \\
    Student Name(s): & \StudentNames \\
    Student ID(s): & \StudentIDs \\
    DVSA Website URL: & \DVSAWebsiteURL \\
    AWS Region: & \AWSRegion \\
    Date: & \SubmissionDate \\
  \end{tabular}

  \vfill
  {\large Project Report\par}
\end{titlepage}

\clearpage
\hypersetup{pageanchor=true}
\pagenumbering{roman}
\tableofcontents
\clearpage

\chapter*{Project Overview}
\addcontentsline{toc}{chapter}{Project Overview}

\section*{Scope}
\placeholder{Summarize the project scope: deployment in a non-production AWS account, the 10 official lessons, remediation, and verification.}

\section*{Deployment Summary}
\begin{itemize}
  \item DVSA Website URL: \DVSAWebsiteURL
  \item AWS region: \AWSRegion
  \item CloudFormation stack: \placeholder{Stack name}
  \item Main services used: \placeholder{S3, API Gateway, Lambda, DynamoDB, SQS, Cognito, IAM, CloudWatch}
\end{itemize}

\section*{Architecture and Request Pipeline}
\placeholder{Explain the intended high-level request flow from browser to API Gateway to Lambda to backend services and back.}
\evidencebox{Optional architecture figure or request-flow diagram.}

\section*{Methodology and Evidence Sources}
\begin{itemize}
  \item \placeholder{How the environment was validated before testing.}
  \item \placeholder{What tools were used to reproduce and verify vulnerabilities.}
  \item \placeholder{How evidence was captured and redacted.}
  \item \placeholder{How fixes were verified after remediation.}
\end{itemize}

\clearpage
\pagenumbering{arabic}

\lessonchapter{01}{Event Injection}{unsafe input handling}
\lessonchapter{02}{Broken Authentication}{authentication failure}
\lessonchapter{03}{Sensitive Information Disclosure}{sensitive data exposure}
\lessonchapter{04}{Insecure Cloud Configuration}{cloud misconfiguration}
\lessonchapter{05}{Broken Access Control}{authorization failure}
\lessonchapter{06}{Denial of Service}{availability failure}
\lessonchapter{07}{Over-Privileged Function}{excessive IAM privilege}
\lessonchapter{08}{Logic Vulnerabilities}{workflow or race-condition failure}
\lessonchapter{09}{Vulnerable Dependencies}{dependency risk}
\lessonchapter{10}{Unhandled Exceptions}{unsafe error handling}

\chapter*{Optional Bonus Findings}
\addcontentsline{toc}{chapter}{Optional Bonus Findings}
\placeholder{Use this section only if you discover valid additional findings beyond the 10 official lessons. Repeat the same reporting discipline if you include bonus work.}

\chapter*{Conclusion}
\addcontentsline{toc}{chapter}{Conclusion}
\placeholder{Summarize the project outcomes, the most important remediation themes, and the main security lessons learned from DVSA.}

\end{document}
